\documentclass[11pt,a4paper]{article}
\usepackage[utf8]{inputenc}
\usepackage{graphicx}
\usepackage{booktabs}
\usepackage{amsmath}
\usepackage{amssymb}
\usepackage{geometry}
\geometry{margin=1in}
\usepackage{float}
\usepackage{hyperref}
\usepackage{xcolor}


\hypersetup{
    colorlinks=true,
    linkcolor=blue,
    filecolor=magenta,      
    urlcolor=cyan,
    pdftitle={Comparative Analysis of Stochastic Soil Water Model Calibration},
    pdfpagemode=FullScreen,
}

\title{\textbf{Stochastic Soil Water Model Calibration across Australian Ecosystems:\\ Empirical Best Fit vs. Eco-Evolutionary Optimality}}
\author{\textbf{Sanjay S.} \\ Department of Civil Engineering, Monash University, Australia \\ \texttt{sanjays@monash.edu}}
\date{\today}

\begin{document}

\maketitle

\begin{abstract}
Predicting plant water-use strategies under water-limited conditions is critical for understanding ecosystem resilience to climate volatility. This report compares two calibration methodologies for a parsimonious Stochastic Soil Water Model (SSWM) across 17 Australian OzFlux sites spanning diverse climates and vegetation types. The first methodology, \textbf{Empirical Best Fit (BF)}, calibrates soil and vegetation traits by maximizing the log-likelihood of observed soil moisture probability density functions (PDFs). The second methodology, \textbf{Eco-Evolutionary Optimality (OPT)}, optimizes parameters to maximize plant transpiration efficiency ($\epsilon$). Our results reveal a systematic divergence between actual plant behavior and theoretical optimality: actual plants exhibit a significantly higher Plant Water Flux Control ($\Pi_F$) than predicted (approximately +1.0 in dimensionless units), translating into a ~74\% reduction in maximum hydraulic capacity. This represents a long-term ``resilience tax'' or insurance policy against hydro-climatic variability (specifically ENSO cycles) that is not accounted for in standard optimality models.
\end{abstract}

\section{Introduction}
Eco-Evolutionary Optimality (EEO) theory posits that plants adjust their leaf stomatal conductance and hydraulic traits to maximize carbon gain per unit water lost, while avoiding the catastrophic risk of hydraulic failure. To test these theories, Bassiouni et al. (2023) developed the One-Water-Universal-Solver (OWUS) utilizing a parsimonious Stochastic Soil Water Model (SSWM). The model simulates soil moisture probability density functions (PDFs) based on random rainfall pulses and plant water extraction curves defined by physiological thresholds:
\begin{itemize}
    \item \textbf{Wilting point ($s_{wilt}$)}: The soil moisture level where plant transpiration ceases.
    \item \textbf{Stress threshold ($s^*$)}: The soil moisture level below which plants begin to regulate stomata and reduce transpiration.
    \item \textbf{Field capacity ($s_{fc}$)}: The soil moisture level above which gravity drains water.
    \item \textbf{Hygroscopic point ($s_h$)}: The absolute lower limit of soil water content.
\end{itemize}

Using Markov Chain Monte Carlo (MCMC) optimization over 17 Australian OzFlux eddy-covariance sites, we calibrated the key parameters: xylem water potential at 50\% loss of conductivity ($P_{x50}$), stomatal closure potential at 50\% reduction in conductance ($P_{g50}$), maximum xylem conductivity ($k_{xl,max}$), and root area index ($RAI$). We compare two alternative calibration frameworks:
\begin{enumerate}
    \item \textbf{Benchmark Empirical Best Fit (BF)}: Parameters are calibrated to directly match the observed soil moisture distribution.
    \item \textbf{Theoretical Optimality (OPT)}: Parameters are optimized to maximize the dimensionless optimality parameter $\epsilon = (1-\theta)E_t/P$, representing unstressed transpiration efficiency.
\end{enumerate}

\section{Methodology and Parameter Definitions}
The model defines several dimensionless ratios that encapsulate water-use strategies:
\begin{equation}
    \Pi_R = \frac{P_{g50}}{P_{x50}}
\end{equation}
representing the ratio of stomatal sensitivity to hydraulic capacity.
\begin{equation}
    \Pi_F = -\frac{E_0}{K_{p,max} P_{g50}}
\end{equation}
representing the Plant Water Flux Control. A high $\Pi_F$ implies a strong hydraulic constraint (low maximum conductance capacity relative to atmospheric demand).
\begin{equation}
    \theta = \text{Dynamic Plant Water Stress}
\end{equation}
representing the proportion of time the plant experiences water stress. The ecosystem resilience index is defined as $R = 1-\theta$.

\section{Results and Performance Comparison}
Table~\ref{tab:perf_comp} compares the Nash-Sutcliffe Efficiency (NSE) and model bias for the two configurations. Table~\ref{tab:traits_comp} compares the calibrated plant traits ($s_{wilt}$, $s^*$, $\Pi_F$, $\epsilon$) between BF and OPT.

% INSERT TABLE 1
\begin{table}[H]
\centering
\caption{Model Calibration Performance Comparison (Nash-Sutcliffe Efficiency and Bias) between Empirical Best Fit (Benchmark) and Eco-Evolutionary Optimality (Optimized) configurations.}
\label{tab:perf_comp}
\begin{tabular}{llccccc}
\toprule
Site ID & IGBP & Benchmark NSE & Optimized NSE & $\Delta$ NSE & Benchmark Bias & Optimized Bias \\
\midrule
AU-Adr & GRA & 0.939 & 0.706 & \textbf{-0.233} & 0.033 & -0.155 \\
AU-Alp & ENF & 0.440 & -0.193 & \textbf{-0.633} & 0.036 & -0.059 \\
AU-Ctr & EBF & 0.360 & 0.428 & 0.068 & 0.199 & 0.134 \\
AU-DaP & GRA & 0.992 & 0.622 & \textbf{-0.369} & -0.005 & -0.145 \\
AU-Eme & GRA & 0.985 & 0.960 & \textbf{-0.025} & -0.038 & -0.066 \\
AU-Fog & GRA & 0.675 & -7.003 & \textbf{-7.678} & -0.019 & -0.233 \\
AU-Gat & GRA & 0.904 & 0.790 & \textbf{-0.114} & -0.071 & -0.108 \\
AU-Gre & SAV & 0.919 & 0.827 & \textbf{-0.091} & -0.103 & -0.173 \\
AU-Lit & SAV & 0.949 & 0.871 & \textbf{-0.078} & -0.022 & 0.149 \\
AU-RDF & GRA & 0.863 & 0.854 & \textbf{-0.009} & -0.178 & -0.000 \\
AU-Rgf & CRO & 0.863 & 0.614 & \textbf{-0.249} & -0.058 & -0.233 \\
AU-SiP & WSA & 0.821 & 0.603 & \textbf{-0.218} & -0.124 & -0.250 \\
AU-Stp & GRA & 0.846 & 0.772 & \textbf{-0.074} & 0.180 & -0.130 \\
AU-Wal & EBF & 0.974 & 0.746 & \textbf{-0.228} & -0.001 & -0.035 \\
AU-Whr & WSA & 0.868 & 0.422 & \textbf{-0.447} & -0.053 & 0.200 \\
AU-YarI & GRA & 0.922 & 0.924 & 0.002 & 0.086 & 0.082 \\
AU-Ync & GRA & 0.990 & 0.631 & \textbf{-0.359} & -0.041 & -0.226 \\
\bottomrule
\end{tabular}
\end{table}

% INSERT TABLE 2
\begin{table}[H]
\centering
\caption{Comparison of calibrated ecohydrological traits and optimality parameters: Wilting point ($s_{wilt}$), stress threshold ($s^*$), Plant Water Flux Control ($\Pi_F$), and Transpiration/Carbon Optimality ($\epsilon$).}
\label{tab:traits_comp}
\begin{tabular}{lcccccccc}
\toprule
 & \multicolumn{2}{c}{$s_{wilt}$} & \multicolumn{2}{c}{$s^*$} & \multicolumn{2}{c}{$\Pi_F$} & \multicolumn{2}{c}{$\epsilon$} \\
\cmidrule(lr){2-3} \cmidrule(lr){4-5} \cmidrule(lr){6-7} \cmidrule(lr){8-9}
Site ID & BF & OPT & BF & OPT & BF & OPT & BF & OPT \\
\midrule
AU-Adr & 0.309 & 0.225 & 0.491 & 0.359 & \textbf{1.282} & 0.975 & 0.216 & \textbf{0.285} \\
AU-Alp & 0.400 & 0.342 & 0.592 & 0.532 & \textbf{1.467} & 1.904 & 0.450 & \textbf{0.494} \\
AU-Ctr & 0.340 & 0.280 & 0.526 & 0.433 & \textbf{0.086} & 0.024 & 0.240 & \textbf{0.267} \\
AU-DaP & 0.237 & 0.185 & 0.384 & 0.297 & \textbf{0.980} & 0.360 & 0.168 & \textbf{0.205} \\
AU-Eme & 0.274 & 0.271 & 0.500 & 0.496 & \textbf{2.217} & 1.808 & 0.124 & \textbf{0.133} \\
AU-Fog & 0.499 & 0.431 & 0.703 & 0.593 & \textbf{3.139} & 0.299 & 0.179 & \textbf{0.528} \\
AU-Gat & 0.236 & 0.204 & 0.395 & 0.334 & \textbf{0.893} & 0.029 & 0.163 & \textbf{0.256} \\
AU-Gre & 0.172 & 0.174 & 0.304 & 0.291 & \textbf{4.408} & 3.080 & 0.004 & \textbf{0.044} \\
AU-Lit & 0.183 & 0.176 & 0.318 & 0.310 & \textbf{0.341} & 0.622 & 0.439 & \textbf{0.448} \\
AU-RDF & 0.237 & 0.187 & 0.368 & 0.297 & \textbf{1.815} & 3.252 & 0.003 & \textbf{0.016} \\
AU-Rgf & 0.044 & 0.035 & 0.113 & 0.089 & \textbf{1.258} & 0.052 & 0.437 & \textbf{0.587} \\
AU-SiP & 0.166 & 0.167 & 0.307 & 0.305 & \textbf{0.563} & 0.358 & 0.655 & \textbf{0.668} \\
AU-Stp & 0.087 & 0.027 & 0.309 & 0.107 & \textbf{1.616} & 1.097 & 0.045 & \textbf{0.158} \\
AU-Wal & 0.587 & 0.566 & 0.783 & 0.769 & \textbf{1.076} & 0.699 & 0.715 & \textbf{0.721} \\
AU-Whr & 0.286 & 0.286 & 0.431 & 0.453 & \textbf{0.125} & 1.532 & 0.172 & \textbf{0.351} \\
AU-YarI & 0.013 & 0.013 & 0.057 & 0.056 & \textbf{0.102} & 0.097 & 0.564 & \textbf{0.566} \\
AU-Ync & 0.109 & 0.089 & 0.298 & 0.241 & \textbf{2.341} & 0.928 & 0.267 & \textbf{0.341} \\
\bottomrule
\end{tabular}
\end{table}

\subsection{Performance Divergence}
The Empirical Best Fit (BF) model consistently reproduces observed soil moisture distributions with high accuracy (NSE $\ge 0.85$ for 14 out of 17 sites). However, when calibrating under the Eco-Evolutionary Optimality (OPT) framework, the model performance degrades significantly for several sites (e.g., Alice Springs Mulga AU-Alp drops to $-0.193$, and Fogg Dam AU-Fog collapses to $-7.003$). This performance gap reflects a fundamental mismatch between the theoretical optimal parameters and actual plant observations.

\subsection{The Hydraulic Capacity Bias}
The most striking result is the systematic difference in the Plant Water Flux Control ($\Pi_F$) between BF and OPT. Under OPT, the theoretical optimum $\Pi_F$ is low (averaging $<1.0$), implying that plants should build highly conductive hydraulic pathways (high $K_{p,max}$) to maximize water throughput. In reality, the BF calibration shows that Australian plants operate with much higher $\Pi_F$ (often $>1.5$, and up to $4.4$ at Great Victoria Desert AU-Gre). This indicates that plants are operating with a \textbf{74\% lower hydraulic capacity} than the theoretical optimum.

\begin{figure}[H]
    \centering
    \includegraphics[width=0.7\textwidth]{01_Main_Paper_Figures/3_piRpiF_range.png}
    \caption{Empirical results of $\Pi_R$ vs $\Pi_F$ across the sites showing the systematic deviation from the global optimal curve (solid black line), indicating a strong conservative hydraulic constraint in Australian vegetation.}
    \label{fig:piRpiF}
\end{figure}

\section{Scientific Discussion}

\subsection{The ``Resilience Tax'' Paradox}
Why do Australian plants underperform relative to the EEO optimum? The optimality model assumes a plant operates to maximize short-term growth and transpiration. In volatile climates like Australia, dominated by multi-year ENSO cycles (decades of dry El Ni\~no and wet La Ni\~na phases), a plant optimized for the average year would perish during severe droughts. 
By maintaining a lower hydraulic capacity ($\Pi_F \gg \Pi_{F,opt}$), plants accept a carbon productivity penalty (operating at ~91\% of maximum potential efficiency, a ~9\% performance gap) to purchase a safety buffer. They pay a \textbf{resilience tax} in the form of conservative xylem construction.

\subsection{Saver vs. Spender Strategies}
Ecosystem hydraulic response falls into two major life-history categories:
\begin{enumerate}
    \item \textbf{Woody Forests as ``Savers''}: Sites like Wallaby Creek (AU-Wal) and Fogg Dam (AU-Fog) invest heavily in carbon-expensive xylem tissue. Cavitation is fatal or extremely slow to repair. They exhibit high $\Pi_F$ to preserve their structural investment, hunkering down during droughts.
    \item \textbf{Grasslands as ``Spenders''}: Sites like Sturt Plains (AU-Stp) and Yanco (AU-Ync) have low structural investment. They can senesce and recover rapidly from seed banks or rhizomes. They operate with lower $\Pi_F$, spending water aggressively when available.
\end{enumerate}

\begin{figure}[H]
    \centering
    \includegraphics[width=0.7\textwidth]{01_Main_Paper_Figures/5_flux_scatter.png}
    \caption{Flux validation scatter comparing modeled vs observed latent heat fluxes across different ecosystems, highlighting the performance of the SSWM model under typical climate regimes.}
    \label{fig:flux_scatter}
\end{figure}

\section{Conclusion and EEO Refinements}
We propose three key advancements for next-generation Eco-Evolutionary Optimality models:
\begin{enumerate}
    \item \textbf{Carbon Cost of Safety}: Include a metabolic construction cost term for building xylem ($\text{Cost}(K_{p,max})$) in the plant objective function.
    \item \textbf{Rhizosphere Soil Constraints}: Decouple soil-root interface resistance ($\Pi_S$) from xylem resistance ($\Pi_F$) in ancient, weathered soils.
    \item \textbf{Dynamic Temporal Plasticity}: Replace static parameters with time-varying parameters that respond to multi-year climate oscillations (ENSO/IOD).
\end{enumerate}

\newpage
\appendix
\section{Individual Site Calibration and Diagnostic Catalog}
This appendix presents the comparison of the soil moisture probability density function (PDF) calibration and Markov Chain Monte Carlo (MCMC) parameter estimation diagnostics for all 17 selected sites. For each site, the diagnostic figures are presented in a top-and-bottom format to ensure maximum readability of the multi-panel grids. The top panel represents the Empirical Best Fit (Benchmark, BF) configuration and the bottom panel represents the Eco-Evolutionary Optimality (Optimized, OPT) configuration. Each site is structured as follows:
\begin{enumerate}
  \item \textbf{Soil Moisture PDF Fits}: Comparison of observed soil moisture histograms and model-fitted PDFs with physiological thresholds ($S_{wilt}$, $S^*$, $S_h$, $S_{fc}$).
  \item \textbf{MCMC Diagnostics (Derived Parameters)}: Prior/posterior parameter distributions, likelihood profiles, and trace chains for derived parameters ($\Pi_T$, $\Pi_S$, $\Pi_R$, $\Pi_F$, $\epsilon$, and $AA$).
  \item \textbf{MCMC Diagnostics (Calibrated Traits)}: Trace diagnostics and distributions for calibrated traits ($P_{x50}$, $P_{g50}$, $k_{xl,max}$, $RAI$, $\beta_{ww}$, $s_{wilt}$, $s^*$, $\psi_{wilt}$, and $\psi^*$).
\end{enumerate}
\newpage
\subsection{AU-Adr --- Adelaide River (GRA)}
Open forest/savanna site located in the Northern Territory, experiencing a distinct monsoonal wet-dry climate. Vegetation is dominated by grassland with scattered eucalyptus trees.
\vskip 1.5em
\begin{figure}[H]
\centering
\includegraphics[width=0.75\textwidth]{output_corrected/bf/figs/_ozflux_1/AU-Adr_GRA_ps.png}
\caption{Empirical Best Fit (BF) soil moisture PDF calibration for AU-Adr.}
\label{fig:ps_bf_AU-Adr}
\end{figure}
\begin{figure}[H]
\centering
\includegraphics[width=0.75\textwidth]{output_corrected/opt/figs/_ozflux_1/AU-Adr_GRA_ps.png}
\caption{Eco-Evolutionary Optimality (OPT) soil moisture PDF calibration for AU-Adr.}
\label{fig:ps_opt_AU-Adr}
\end{figure}
\newpage
\begin{figure}[H]
\centering
\includegraphics[width=0.9\textwidth]{output_corrected/bf/figs/_ozflux_1/AU-Adr_GRA_1.png}
\caption{MCMC diagnostics for derived parameters ($\Pi_T$, $\Pi_S$, $\Pi_R$, $\Pi_F$, $\epsilon$, and $AA$) under Empirical Best Fit (BF) for AU-Adr.}
\label{fig:1_bf_AU-Adr}
\end{figure}
\begin{figure}[H]
\centering
\includegraphics[width=0.9\textwidth]{output_corrected/opt/figs/_ozflux_1/AU-Adr_GRA_1.png}
\caption{MCMC diagnostics for derived parameters under Eco-Evolutionary Optimality (OPT) for AU-Adr.}
\label{fig:1_opt_AU-Adr}
\end{figure}
\newpage
\begin{figure}[H]
\centering
\includegraphics[width=0.9\textwidth]{output_corrected/bf/figs/_ozflux_1/AU-Adr_GRA_2.png}
\caption{MCMC diagnostics for calibrated traits ($P_{x50}$, $P_{g50}$, $k_{xl,max}$, $RAI$, $\beta_{ww}$, $s_{wilt}$, $s^*$, $\psi_{wilt}$, and $\psi^*$) under Empirical Best Fit (BF) for AU-Adr.}
\label{fig:2_bf_AU-Adr}
\end{figure}
\begin{figure}[H]
\centering
\includegraphics[width=0.9\textwidth]{output_corrected/opt/figs/_ozflux_1/AU-Adr_GRA_2.png}
\caption{MCMC diagnostics for calibrated traits under Eco-Evolutionary Optimality (OPT) for AU-Adr.}
\label{fig:2_opt_AU-Adr}
\end{figure}
\newpage
\subsection{AU-Alp --- Alice Springs Mulga (ENF)}
Arid grassland and Mulga (Acacia aneura) woodland site in central Northern Territory. Characterized by low, highly episodic rainfall and extreme temperatures.
\vskip 1.5em
\begin{figure}[H]
\centering
\includegraphics[width=0.75\textwidth]{output_corrected/bf/figs/_ozflux_1/AU-Alp_ENF_ps.png}
\caption{Empirical Best Fit (BF) soil moisture PDF calibration for AU-Alp.}
\label{fig:ps_bf_AU-Alp}
\end{figure}
\begin{figure}[H]
\centering
\includegraphics[width=0.75\textwidth]{output_corrected/opt/figs/_ozflux_1/AU-Alp_ENF_ps.png}
\caption{Eco-Evolutionary Optimality (OPT) soil moisture PDF calibration for AU-Alp.}
\label{fig:ps_opt_AU-Alp}
\end{figure}
\newpage
\begin{figure}[H]
\centering
\includegraphics[width=0.9\textwidth]{output_corrected/bf/figs/_ozflux_1/AU-Alp_ENF_1.png}
\caption{MCMC diagnostics for derived parameters ($\Pi_T$, $\Pi_S$, $\Pi_R$, $\Pi_F$, $\epsilon$, and $AA$) under Empirical Best Fit (BF) for AU-Alp.}
\label{fig:1_bf_AU-Alp}
\end{figure}
\begin{figure}[H]
\centering
\includegraphics[width=0.9\textwidth]{output_corrected/opt/figs/_ozflux_1/AU-Alp_ENF_1.png}
\caption{MCMC diagnostics for derived parameters under Eco-Evolutionary Optimality (OPT) for AU-Alp.}
\label{fig:1_opt_AU-Alp}
\end{figure}
\newpage
\begin{figure}[H]
\centering
\includegraphics[width=0.9\textwidth]{output_corrected/bf/figs/_ozflux_1/AU-Alp_ENF_2.png}
\caption{MCMC diagnostics for calibrated traits ($P_{x50}$, $P_{g50}$, $k_{xl,max}$, $RAI$, $\beta_{ww}$, $s_{wilt}$, $s^*$, $\psi_{wilt}$, and $\psi^*$) under Empirical Best Fit (BF) for AU-Alp.}
\label{fig:2_bf_AU-Alp}
\end{figure}
\begin{figure}[H]
\centering
\includegraphics[width=0.9\textwidth]{output_corrected/opt/figs/_ozflux_1/AU-Alp_ENF_2.png}
\caption{MCMC diagnostics for calibrated traits under Eco-Evolutionary Optimality (OPT) for AU-Alp.}
\label{fig:2_opt_AU-Alp}
\end{figure}
\newpage
\subsection{AU-Ctr --- Calperum Chowilla (EBF)}
Mallee woodland and shrubland site in South Australia. Located in a semi-arid zone, the site has highly variable rainfall and coarse, sandy soils.
\vskip 1.5em
\begin{figure}[H]
\centering
\includegraphics[width=0.75\textwidth]{output_corrected/bf/figs/_ozflux_1/AU-Ctr_EBF_ps.png}
\caption{Empirical Best Fit (BF) soil moisture PDF calibration for AU-Ctr.}
\label{fig:ps_bf_AU-Ctr}
\end{figure}
\begin{figure}[H]
\centering
\includegraphics[width=0.75\textwidth]{output_corrected/opt/figs/_ozflux_1/AU-Ctr_EBF_ps.png}
\caption{Eco-Evolutionary Optimality (OPT) soil moisture PDF calibration for AU-Ctr.}
\label{fig:ps_opt_AU-Ctr}
\end{figure}
\newpage
\begin{figure}[H]
\centering
\includegraphics[width=0.9\textwidth]{output_corrected/bf/figs/_ozflux_1/AU-Ctr_EBF_1.png}
\caption{MCMC diagnostics for derived parameters ($\Pi_T$, $\Pi_S$, $\Pi_R$, $\Pi_F$, $\epsilon$, and $AA$) under Empirical Best Fit (BF) for AU-Ctr.}
\label{fig:1_bf_AU-Ctr}
\end{figure}
\begin{figure}[H]
\centering
\includegraphics[width=0.9\textwidth]{output_corrected/opt/figs/_ozflux_1/AU-Ctr_EBF_1.png}
\caption{MCMC diagnostics for derived parameters under Eco-Evolutionary Optimality (OPT) for AU-Ctr.}
\label{fig:1_opt_AU-Ctr}
\end{figure}
\newpage
\begin{figure}[H]
\centering
\includegraphics[width=0.9\textwidth]{output_corrected/bf/figs/_ozflux_1/AU-Ctr_EBF_2.png}
\caption{MCMC diagnostics for calibrated traits ($P_{x50}$, $P_{g50}$, $k_{xl,max}$, $RAI$, $\beta_{ww}$, $s_{wilt}$, $s^*$, $\psi_{wilt}$, and $\psi^*$) under Empirical Best Fit (BF) for AU-Ctr.}
\label{fig:2_bf_AU-Ctr}
\end{figure}
\begin{figure}[H]
\centering
\includegraphics[width=0.9\textwidth]{output_corrected/opt/figs/_ozflux_1/AU-Ctr_EBF_2.png}
\caption{MCMC diagnostics for calibrated traits under Eco-Evolutionary Optimality (OPT) for AU-Ctr.}
\label{fig:2_opt_AU-Ctr}
\end{figure}
\newpage
\subsection{AU-DaP --- Daly River Cleared (GRA)}
Tropical pasture/grassland site in the Northern Territory. Cleared from native savanna, it undergoes intense monsoonal wet seasons followed by severe dry periods.
\vskip 1.5em
\begin{figure}[H]
\centering
\includegraphics[width=0.75\textwidth]{output_corrected/bf/figs/_ozflux_1/AU-DaP_GRA_ps.png}
\caption{Empirical Best Fit (BF) soil moisture PDF calibration for AU-DaP.}
\label{fig:ps_bf_AU-DaP}
\end{figure}
\begin{figure}[H]
\centering
\includegraphics[width=0.75\textwidth]{output_corrected/opt/figs/_ozflux_1/AU-DaP_GRA_ps.png}
\caption{Eco-Evolutionary Optimality (OPT) soil moisture PDF calibration for AU-DaP.}
\label{fig:ps_opt_AU-DaP}
\end{figure}
\newpage
\begin{figure}[H]
\centering
\includegraphics[width=0.9\textwidth]{output_corrected/bf/figs/_ozflux_1/AU-DaP_GRA_1.png}
\caption{MCMC diagnostics for derived parameters ($\Pi_T$, $\Pi_S$, $\Pi_R$, $\Pi_F$, $\epsilon$, and $AA$) under Empirical Best Fit (BF) for AU-DaP.}
\label{fig:1_bf_AU-DaP}
\end{figure}
\begin{figure}[H]
\centering
\includegraphics[width=0.9\textwidth]{output_corrected/opt/figs/_ozflux_1/AU-DaP_GRA_1.png}
\caption{MCMC diagnostics for derived parameters under Eco-Evolutionary Optimality (OPT) for AU-DaP.}
\label{fig:1_opt_AU-DaP}
\end{figure}
\newpage
\begin{figure}[H]
\centering
\includegraphics[width=0.9\textwidth]{output_corrected/bf/figs/_ozflux_1/AU-DaP_GRA_2.png}
\caption{MCMC diagnostics for calibrated traits ($P_{x50}$, $P_{g50}$, $k_{xl,max}$, $RAI$, $\beta_{ww}$, $s_{wilt}$, $s^*$, $\psi_{wilt}$, and $\psi^*$) under Empirical Best Fit (BF) for AU-DaP.}
\label{fig:2_bf_AU-DaP}
\end{figure}
\begin{figure}[H]
\centering
\includegraphics[width=0.9\textwidth]{output_corrected/opt/figs/_ozflux_1/AU-DaP_GRA_2.png}
\caption{MCMC diagnostics for calibrated traits under Eco-Evolutionary Optimality (OPT) for AU-DaP.}
\label{fig:2_opt_AU-DaP}
\end{figure}
\newpage
\subsection{AU-Eme --- Emerald (GRA)}
Dry tropical agricultural grassland/pasture site in central Queensland. Characterized by high solar radiation and summer-dominant rainfall.
\vskip 1.5em
\begin{figure}[H]
\centering
\includegraphics[width=0.75\textwidth]{output_corrected/bf/figs/_ozflux_1/AU-Eme_GRA_ps.png}
\caption{Empirical Best Fit (BF) soil moisture PDF calibration for AU-Eme.}
\label{fig:ps_bf_AU-Eme}
\end{figure}
\begin{figure}[H]
\centering
\includegraphics[width=0.75\textwidth]{output_corrected/opt/figs/_ozflux_1/AU-Eme_GRA_ps.png}
\caption{Eco-Evolutionary Optimality (OPT) soil moisture PDF calibration for AU-Eme.}
\label{fig:ps_opt_AU-Eme}
\end{figure}
\newpage
\begin{figure}[H]
\centering
\includegraphics[width=0.9\textwidth]{output_corrected/bf/figs/_ozflux_1/AU-Eme_GRA_1.png}
\caption{MCMC diagnostics for derived parameters ($\Pi_T$, $\Pi_S$, $\Pi_R$, $\Pi_F$, $\epsilon$, and $AA$) under Empirical Best Fit (BF) for AU-Eme.}
\label{fig:1_bf_AU-Eme}
\end{figure}
\begin{figure}[H]
\centering
\includegraphics[width=0.9\textwidth]{output_corrected/opt/figs/_ozflux_1/AU-Eme_GRA_1.png}
\caption{MCMC diagnostics for derived parameters under Eco-Evolutionary Optimality (OPT) for AU-Eme.}
\label{fig:1_opt_AU-Eme}
\end{figure}
\newpage
\begin{figure}[H]
\centering
\includegraphics[width=0.9\textwidth]{output_corrected/bf/figs/_ozflux_1/AU-Eme_GRA_2.png}
\caption{MCMC diagnostics for calibrated traits ($P_{x50}$, $P_{g50}$, $k_{xl,max}$, $RAI$, $\beta_{ww}$, $s_{wilt}$, $s^*$, $\psi_{wilt}$, and $\psi^*$) under Empirical Best Fit (BF) for AU-Eme.}
\label{fig:2_bf_AU-Eme}
\end{figure}
\begin{figure}[H]
\centering
\includegraphics[width=0.9\textwidth]{output_corrected/opt/figs/_ozflux_1/AU-Eme_GRA_2.png}
\caption{MCMC diagnostics for calibrated traits under Eco-Evolutionary Optimality (OPT) for AU-Eme.}
\label{fig:2_opt_AU-Eme}
\end{figure}
\newpage
\subsection{AU-Fog --- Fogg Dam (GRA)}
Humid tropical flood plain and wetland site in the Northern Territory. Dominated by dense grasslands/forests that remain wet for long periods of the year.
\vskip 1.5em
\begin{figure}[H]
\centering
\includegraphics[width=0.75\textwidth]{output_corrected/bf/figs/_ozflux_1/AU-Fog_GRA_ps.png}
\caption{Empirical Best Fit (BF) soil moisture PDF calibration for AU-Fog.}
\label{fig:ps_bf_AU-Fog}
\end{figure}
\begin{figure}[H]
\centering
\includegraphics[width=0.75\textwidth]{output_corrected/opt/figs/_ozflux_1/AU-Fog_GRA_ps.png}
\caption{Eco-Evolutionary Optimality (OPT) soil moisture PDF calibration for AU-Fog.}
\label{fig:ps_opt_AU-Fog}
\end{figure}
\newpage
\begin{figure}[H]
\centering
\includegraphics[width=0.9\textwidth]{output_corrected/bf/figs/_ozflux_1/AU-Fog_GRA_1.png}
\caption{MCMC diagnostics for derived parameters ($\Pi_T$, $\Pi_S$, $\Pi_R$, $\Pi_F$, $\epsilon$, and $AA$) under Empirical Best Fit (BF) for AU-Fog.}
\label{fig:1_bf_AU-Fog}
\end{figure}
\begin{figure}[H]
\centering
\includegraphics[width=0.9\textwidth]{output_corrected/opt/figs/_ozflux_1/AU-Fog_GRA_1.png}
\caption{MCMC diagnostics for derived parameters under Eco-Evolutionary Optimality (OPT) for AU-Fog.}
\label{fig:1_opt_AU-Fog}
\end{figure}
\newpage
\begin{figure}[H]
\centering
\includegraphics[width=0.9\textwidth]{output_corrected/bf/figs/_ozflux_1/AU-Fog_GRA_2.png}
\caption{MCMC diagnostics for calibrated traits ($P_{x50}$, $P_{g50}$, $k_{xl,max}$, $RAI$, $\beta_{ww}$, $s_{wilt}$, $s^*$, $\psi_{wilt}$, and $\psi^*$) under Empirical Best Fit (BF) for AU-Fog.}
\label{fig:2_bf_AU-Fog}
\end{figure}
\begin{figure}[H]
\centering
\includegraphics[width=0.9\textwidth]{output_corrected/opt/figs/_ozflux_1/AU-Fog_GRA_2.png}
\caption{MCMC diagnostics for calibrated traits under Eco-Evolutionary Optimality (OPT) for AU-Fog.}
\label{fig:2_opt_AU-Fog}
\end{figure}
\newpage
\subsection{AU-Gat --- Gatton (GRA)}
Agricultural crop and pasture site in Queensland, situated in a sub-tropical climate zone with clay soils.
\vskip 1.5em
\begin{figure}[H]
\centering
\includegraphics[width=0.75\textwidth]{output_corrected/bf/figs/_ozflux_1/AU-Gat_GRA_ps.png}
\caption{Empirical Best Fit (BF) soil moisture PDF calibration for AU-Gat.}
\label{fig:ps_bf_AU-Gat}
\end{figure}
\begin{figure}[H]
\centering
\includegraphics[width=0.75\textwidth]{output_corrected/opt/figs/_ozflux_1/AU-Gat_GRA_ps.png}
\caption{Eco-Evolutionary Optimality (OPT) soil moisture PDF calibration for AU-Gat.}
\label{fig:ps_opt_AU-Gat}
\end{figure}
\newpage
\begin{figure}[H]
\centering
\includegraphics[width=0.9\textwidth]{output_corrected/bf/figs/_ozflux_1/AU-Gat_GRA_1.png}
\caption{MCMC diagnostics for derived parameters ($\Pi_T$, $\Pi_S$, $\Pi_R$, $\Pi_F$, $\epsilon$, and $AA$) under Empirical Best Fit (BF) for AU-Gat.}
\label{fig:1_bf_AU-Gat}
\end{figure}
\begin{figure}[H]
\centering
\includegraphics[width=0.9\textwidth]{output_corrected/opt/figs/_ozflux_1/AU-Gat_GRA_1.png}
\caption{MCMC diagnostics for derived parameters under Eco-Evolutionary Optimality (OPT) for AU-Gat.}
\label{fig:1_opt_AU-Gat}
\end{figure}
\newpage
\begin{figure}[H]
\centering
\includegraphics[width=0.9\textwidth]{output_corrected/bf/figs/_ozflux_1/AU-Gat_GRA_2.png}
\caption{MCMC diagnostics for calibrated traits ($P_{x50}$, $P_{g50}$, $k_{xl,max}$, $RAI$, $\beta_{ww}$, $s_{wilt}$, $s^*$, $\psi_{wilt}$, and $\psi^*$) under Empirical Best Fit (BF) for AU-Gat.}
\label{fig:2_bf_AU-Gat}
\end{figure}
\begin{figure}[H]
\centering
\includegraphics[width=0.9\textwidth]{output_corrected/opt/figs/_ozflux_1/AU-Gat_GRA_2.png}
\caption{MCMC diagnostics for calibrated traits under Eco-Evolutionary Optimality (OPT) for AU-Gat.}
\label{fig:2_opt_AU-Gat}
\end{figure}
\newpage
\subsection{AU-Gre --- Great Victoria Desert (SAV)}
Arid savanna site in Western Australia. Characterized by hummock grasslands (Spinifex) and sparse shrubs over sand dunes.
\vskip 1.5em
\begin{figure}[H]
\centering
\includegraphics[width=0.75\textwidth]{output_corrected/bf/figs/_ozflux_1/AU-Gre_SAV_ps.png}
\caption{Empirical Best Fit (BF) soil moisture PDF calibration for AU-Gre.}
\label{fig:ps_bf_AU-Gre}
\end{figure}
\begin{figure}[H]
\centering
\includegraphics[width=0.75\textwidth]{output_corrected/opt/figs/_ozflux_1/AU-Gre_SAV_ps.png}
\caption{Eco-Evolutionary Optimality (OPT) soil moisture PDF calibration for AU-Gre.}
\label{fig:ps_opt_AU-Gre}
\end{figure}
\newpage
\begin{figure}[H]
\centering
\includegraphics[width=0.9\textwidth]{output_corrected/bf/figs/_ozflux_1/AU-Gre_SAV_1.png}
\caption{MCMC diagnostics for derived parameters ($\Pi_T$, $\Pi_S$, $\Pi_R$, $\Pi_F$, $\epsilon$, and $AA$) under Empirical Best Fit (BF) for AU-Gre.}
\label{fig:1_bf_AU-Gre}
\end{figure}
\begin{figure}[H]
\centering
\includegraphics[width=0.9\textwidth]{output_corrected/opt/figs/_ozflux_1/AU-Gre_SAV_1.png}
\caption{MCMC diagnostics for derived parameters under Eco-Evolutionary Optimality (OPT) for AU-Gre.}
\label{fig:1_opt_AU-Gre}
\end{figure}
\newpage
\begin{figure}[H]
\centering
\includegraphics[width=0.9\textwidth]{output_corrected/bf/figs/_ozflux_1/AU-Gre_SAV_2.png}
\caption{MCMC diagnostics for calibrated traits ($P_{x50}$, $P_{g50}$, $k_{xl,max}$, $RAI$, $\beta_{ww}$, $s_{wilt}$, $s^*$, $\psi_{wilt}$, and $\psi^*$) under Empirical Best Fit (BF) for AU-Gre.}
\label{fig:2_bf_AU-Gre}
\end{figure}
\begin{figure}[H]
\centering
\includegraphics[width=0.9\textwidth]{output_corrected/opt/figs/_ozflux_1/AU-Gre_SAV_2.png}
\caption{MCMC diagnostics for calibrated traits under Eco-Evolutionary Optimality (OPT) for AU-Gre.}
\label{fig:2_opt_AU-Gre}
\end{figure}
\newpage
\subsection{AU-Lit --- Litchfield (SAV)}
Tropical savanna site in the Northern Territory, dominated by open eucalyptus forest and tall grass understory, subject to annual fire cycles.
\vskip 1.5em
\begin{figure}[H]
\centering
\includegraphics[width=0.75\textwidth]{output_corrected/bf/figs/_ozflux_1/AU-Lit_SAV_ps.png}
\caption{Empirical Best Fit (BF) soil moisture PDF calibration for AU-Lit.}
\label{fig:ps_bf_AU-Lit}
\end{figure}
\begin{figure}[H]
\centering
\includegraphics[width=0.75\textwidth]{output_corrected/opt/figs/_ozflux_1/AU-Lit_SAV_ps.png}
\caption{Eco-Evolutionary Optimality (OPT) soil moisture PDF calibration for AU-Lit.}
\label{fig:ps_opt_AU-Lit}
\end{figure}
\newpage
\begin{figure}[H]
\centering
\includegraphics[width=0.9\textwidth]{output_corrected/bf/figs/_ozflux_1/AU-Lit_SAV_1.png}
\caption{MCMC diagnostics for derived parameters ($\Pi_T$, $\Pi_S$, $\Pi_R$, $\Pi_F$, $\epsilon$, and $AA$) under Empirical Best Fit (BF) for AU-Lit.}
\label{fig:1_bf_AU-Lit}
\end{figure}
\begin{figure}[H]
\centering
\includegraphics[width=0.9\textwidth]{output_corrected/opt/figs/_ozflux_1/AU-Lit_SAV_1.png}
\caption{MCMC diagnostics for derived parameters under Eco-Evolutionary Optimality (OPT) for AU-Lit.}
\label{fig:1_opt_AU-Lit}
\end{figure}
\newpage
\begin{figure}[H]
\centering
\includegraphics[width=0.9\textwidth]{output_corrected/bf/figs/_ozflux_1/AU-Lit_SAV_2.png}
\caption{MCMC diagnostics for calibrated traits ($P_{x50}$, $P_{g50}$, $k_{xl,max}$, $RAI$, $\beta_{ww}$, $s_{wilt}$, $s^*$, $\psi_{wilt}$, and $\psi^*$) under Empirical Best Fit (BF) for AU-Lit.}
\label{fig:2_bf_AU-Lit}
\end{figure}
\begin{figure}[H]
\centering
\includegraphics[width=0.9\textwidth]{output_corrected/opt/figs/_ozflux_1/AU-Lit_SAV_2.png}
\caption{MCMC diagnostics for calibrated traits under Eco-Evolutionary Optimality (OPT) for AU-Lit.}
\label{fig:2_opt_AU-Lit}
\end{figure}
\newpage
\subsection{AU-RDF --- Red Dirt Melon Farm (GRA)}
Cleared pasture site in the Northern Territory. Highly seasonal monsoonal climate with sandy soils and high evaporative demand.
\vskip 1.5em
\begin{figure}[H]
\centering
\includegraphics[width=0.75\textwidth]{output_corrected/bf/figs/_ozflux_1/AU-RDF_GRA_ps.png}
\caption{Empirical Best Fit (BF) soil moisture PDF calibration for AU-RDF.}
\label{fig:ps_bf_AU-RDF}
\end{figure}
\begin{figure}[H]
\centering
\includegraphics[width=0.75\textwidth]{output_corrected/opt/figs/_ozflux_1/AU-RDF_GRA_ps.png}
\caption{Eco-Evolutionary Optimality (OPT) soil moisture PDF calibration for AU-RDF.}
\label{fig:ps_opt_AU-RDF}
\end{figure}
\newpage
\begin{figure}[H]
\centering
\includegraphics[width=0.9\textwidth]{output_corrected/bf/figs/_ozflux_1/AU-RDF_GRA_1.png}
\caption{MCMC diagnostics for derived parameters ($\Pi_T$, $\Pi_S$, $\Pi_R$, $\Pi_F$, $\epsilon$, and $AA$) under Empirical Best Fit (BF) for AU-RDF.}
\label{fig:1_bf_AU-RDF}
\end{figure}
\begin{figure}[H]
\centering
\includegraphics[width=0.9\textwidth]{output_corrected/opt/figs/_ozflux_1/AU-RDF_GRA_1.png}
\caption{MCMC diagnostics for derived parameters under Eco-Evolutionary Optimality (OPT) for AU-RDF.}
\label{fig:1_opt_AU-RDF}
\end{figure}
\newpage
\begin{figure}[H]
\centering
\includegraphics[width=0.9\textwidth]{output_corrected/bf/figs/_ozflux_1/AU-RDF_GRA_2.png}
\caption{MCMC diagnostics for calibrated traits ($P_{x50}$, $P_{g50}$, $k_{xl,max}$, $RAI$, $\beta_{ww}$, $s_{wilt}$, $s^*$, $\psi_{wilt}$, and $\psi^*$) under Empirical Best Fit (BF) for AU-RDF.}
\label{fig:2_bf_AU-RDF}
\end{figure}
\begin{figure}[H]
\centering
\includegraphics[width=0.9\textwidth]{output_corrected/opt/figs/_ozflux_1/AU-RDF_GRA_2.png}
\caption{MCMC diagnostics for calibrated traits under Eco-Evolutionary Optimality (OPT) for AU-RDF.}
\label{fig:2_opt_AU-RDF}
\end{figure}
\newpage
\subsection{AU-Rgf --- Ridgefield (CRO)}
Mediterranean cropland site in Western Australia. Experiences wet winters and dry summers, dominated by annual crop species.
\vskip 1.5em
\begin{figure}[H]
\centering
\includegraphics[width=0.75\textwidth]{output_corrected/bf/figs/_ozflux_1/AU-Rgf_CRO_ps.png}
\caption{Empirical Best Fit (BF) soil moisture PDF calibration for AU-Rgf.}
\label{fig:ps_bf_AU-Rgf}
\end{figure}
\begin{figure}[H]
\centering
\includegraphics[width=0.75\textwidth]{output_corrected/opt/figs/_ozflux_1/AU-Rgf_CRO_ps.png}
\caption{Eco-Evolutionary Optimality (OPT) soil moisture PDF calibration for AU-Rgf.}
\label{fig:ps_opt_AU-Rgf}
\end{figure}
\newpage
\begin{figure}[H]
\centering
\includegraphics[width=0.9\textwidth]{output_corrected/bf/figs/_ozflux_1/AU-Rgf_CRO_1.png}
\caption{MCMC diagnostics for derived parameters ($\Pi_T$, $\Pi_S$, $\Pi_R$, $\Pi_F$, $\epsilon$, and $AA$) under Empirical Best Fit (BF) for AU-Rgf.}
\label{fig:1_bf_AU-Rgf}
\end{figure}
\begin{figure}[H]
\centering
\includegraphics[width=0.9\textwidth]{output_corrected/opt/figs/_ozflux_1/AU-Rgf_CRO_1.png}
\caption{MCMC diagnostics for derived parameters under Eco-Evolutionary Optimality (OPT) for AU-Rgf.}
\label{fig:1_opt_AU-Rgf}
\end{figure}
\newpage
\begin{figure}[H]
\centering
\includegraphics[width=0.9\textwidth]{output_corrected/bf/figs/_ozflux_1/AU-Rgf_CRO_2.png}
\caption{MCMC diagnostics for calibrated traits ($P_{x50}$, $P_{g50}$, $k_{xl,max}$, $RAI$, $\beta_{ww}$, $s_{wilt}$, $s^*$, $\psi_{wilt}$, and $\psi^*$) under Empirical Best Fit (BF) for AU-Rgf.}
\label{fig:2_bf_AU-Rgf}
\end{figure}
\begin{figure}[H]
\centering
\includegraphics[width=0.9\textwidth]{output_corrected/opt/figs/_ozflux_1/AU-Rgf_CRO_2.png}
\caption{MCMC diagnostics for calibrated traits under Eco-Evolutionary Optimality (OPT) for AU-Rgf.}
\label{fig:2_opt_AU-Rgf}
\end{figure}
\newpage
\subsection{AU-SiP --- Silver Plains (WSA)}
Woody savanna site in Cape York, Queensland. Subject to tropical wet seasons and extensive seasonal drying.
\vskip 1.5em
\begin{figure}[H]
\centering
\includegraphics[width=0.75\textwidth]{output_corrected/bf/figs/_ozflux_1/AU-SiP_WSA_ps.png}
\caption{Empirical Best Fit (BF) soil moisture PDF calibration for AU-SiP.}
\label{fig:ps_bf_AU-SiP}
\end{figure}
\begin{figure}[H]
\centering
\includegraphics[width=0.75\textwidth]{output_corrected/opt/figs/_ozflux_1/AU-SiP_WSA_ps.png}
\caption{Eco-Evolutionary Optimality (OPT) soil moisture PDF calibration for AU-SiP.}
\label{fig:ps_opt_AU-SiP}
\end{figure}
\newpage
\begin{figure}[H]
\centering
\includegraphics[width=0.9\textwidth]{output_corrected/bf/figs/_ozflux_1/AU-SiP_WSA_1.png}
\caption{MCMC diagnostics for derived parameters ($\Pi_T$, $\Pi_S$, $\Pi_R$, $\Pi_F$, $\epsilon$, and $AA$) under Empirical Best Fit (BF) for AU-SiP.}
\label{fig:1_bf_AU-SiP}
\end{figure}
\begin{figure}[H]
\centering
\includegraphics[width=0.9\textwidth]{output_corrected/opt/figs/_ozflux_1/AU-SiP_WSA_1.png}
\caption{MCMC diagnostics for derived parameters under Eco-Evolutionary Optimality (OPT) for AU-SiP.}
\label{fig:1_opt_AU-SiP}
\end{figure}
\newpage
\begin{figure}[H]
\centering
\includegraphics[width=0.9\textwidth]{output_corrected/bf/figs/_ozflux_1/AU-SiP_WSA_2.png}
\caption{MCMC diagnostics for calibrated traits ($P_{x50}$, $P_{g50}$, $k_{xl,max}$, $RAI$, $\beta_{ww}$, $s_{wilt}$, $s^*$, $\psi_{wilt}$, and $\psi^*$) under Empirical Best Fit (BF) for AU-SiP.}
\label{fig:2_bf_AU-SiP}
\end{figure}
\begin{figure}[H]
\centering
\includegraphics[width=0.9\textwidth]{output_corrected/opt/figs/_ozflux_1/AU-SiP_WSA_2.png}
\caption{MCMC diagnostics for calibrated traits under Eco-Evolutionary Optimality (OPT) for AU-SiP.}
\label{fig:2_opt_AU-SiP}
\end{figure}
\newpage
\subsection{AU-Stp --- Sturt Plains (GRA)}
Cracking clay grassland site in the Northern Territory. Experiences extreme seasonal hydrology (flooding to deep cracking dry-down).
\vskip 1.5em
\begin{figure}[H]
\centering
\includegraphics[width=0.75\textwidth]{output_corrected/bf/figs/_ozflux_1/AU-Stp_GRA_ps.png}
\caption{Empirical Best Fit (BF) soil moisture PDF calibration for AU-Stp.}
\label{fig:ps_bf_AU-Stp}
\end{figure}
\begin{figure}[H]
\centering
\includegraphics[width=0.75\textwidth]{output_corrected/opt/figs/_ozflux_1/AU-Stp_GRA_ps.png}
\caption{Eco-Evolutionary Optimality (OPT) soil moisture PDF calibration for AU-Stp.}
\label{fig:ps_opt_AU-Stp}
\end{figure}
\newpage
\begin{figure}[H]
\centering
\includegraphics[width=0.9\textwidth]{output_corrected/bf/figs/_ozflux_1/AU-Stp_GRA_1.png}
\caption{MCMC diagnostics for derived parameters ($\Pi_T$, $\Pi_S$, $\Pi_R$, $\Pi_F$, $\epsilon$, and $AA$) under Empirical Best Fit (BF) for AU-Stp.}
\label{fig:1_bf_AU-Stp}
\end{figure}
\begin{figure}[H]
\centering
\includegraphics[width=0.9\textwidth]{output_corrected/opt/figs/_ozflux_1/AU-Stp_GRA_1.png}
\caption{MCMC diagnostics for derived parameters under Eco-Evolutionary Optimality (OPT) for AU-Stp.}
\label{fig:1_opt_AU-Stp}
\end{figure}
\newpage
\begin{figure}[H]
\centering
\includegraphics[width=0.9\textwidth]{output_corrected/bf/figs/_ozflux_1/AU-Stp_GRA_2.png}
\caption{MCMC diagnostics for calibrated traits ($P_{x50}$, $P_{g50}$, $k_{xl,max}$, $RAI$, $\beta_{ww}$, $s_{wilt}$, $s^*$, $\psi_{wilt}$, and $\psi^*$) under Empirical Best Fit (BF) for AU-Stp.}
\label{fig:2_bf_AU-Stp}
\end{figure}
\begin{figure}[H]
\centering
\includegraphics[width=0.9\textwidth]{output_corrected/opt/figs/_ozflux_1/AU-Stp_GRA_2.png}
\caption{MCMC diagnostics for calibrated traits under Eco-Evolutionary Optimality (OPT) for AU-Stp.}
\label{fig:2_opt_AU-Stp}
\end{figure}
\newpage
\subsection{AU-Wal --- Wallaby Creek (EBF)}
Wet sclerophyll eucalyptus forest in Victoria. Located in a temperate mountainous region with high rainfall and deep clay loam soils.
\vskip 1.5em
\begin{figure}[H]
\centering
\includegraphics[width=0.75\textwidth]{output_corrected/bf/figs/_ozflux_1/AU-Wal_EBF_ps.png}
\caption{Empirical Best Fit (BF) soil moisture PDF calibration for AU-Wal.}
\label{fig:ps_bf_AU-Wal}
\end{figure}
\begin{figure}[H]
\centering
\includegraphics[width=0.75\textwidth]{output_corrected/opt/figs/_ozflux_1/AU-Wal_EBF_ps.png}
\caption{Eco-Evolutionary Optimality (OPT) soil moisture PDF calibration for AU-Wal.}
\label{fig:ps_opt_AU-Wal}
\end{figure}
\newpage
\begin{figure}[H]
\centering
\includegraphics[width=0.9\textwidth]{output_corrected/bf/figs/_ozflux_1/AU-Wal_EBF_1.png}
\caption{MCMC diagnostics for derived parameters ($\Pi_T$, $\Pi_S$, $\Pi_R$, $\Pi_F$, $\epsilon$, and $AA$) under Empirical Best Fit (BF) for AU-Wal.}
\label{fig:1_bf_AU-Wal}
\end{figure}
\begin{figure}[H]
\centering
\includegraphics[width=0.9\textwidth]{output_corrected/opt/figs/_ozflux_1/AU-Wal_EBF_1.png}
\caption{MCMC diagnostics for derived parameters under Eco-Evolutionary Optimality (OPT) for AU-Wal.}
\label{fig:1_opt_AU-Wal}
\end{figure}
\newpage
\begin{figure}[H]
\centering
\includegraphics[width=0.9\textwidth]{output_corrected/bf/figs/_ozflux_1/AU-Wal_EBF_2.png}
\caption{MCMC diagnostics for calibrated traits ($P_{x50}$, $P_{g50}$, $k_{xl,max}$, $RAI$, $\beta_{ww}$, $s_{wilt}$, $s^*$, $\psi_{wilt}$, and $\psi^*$) under Empirical Best Fit (BF) for AU-Wal.}
\label{fig:2_bf_AU-Wal}
\end{figure}
\begin{figure}[H]
\centering
\includegraphics[width=0.9\textwidth]{output_corrected/opt/figs/_ozflux_1/AU-Wal_EBF_2.png}
\caption{MCMC diagnostics for calibrated traits under Eco-Evolutionary Optimality (OPT) for AU-Wal.}
\label{fig:2_opt_AU-Wal}
\end{figure}
\newpage
\subsection{AU-Whr --- Whroo (WSA)}
Dry temperate eucalyptus woodland site in Victoria. Sandy clay loam soils with winter-dominant rainfall and hot, dry summers.
\vskip 1.5em
\begin{figure}[H]
\centering
\includegraphics[width=0.75\textwidth]{output_corrected/bf/figs/_ozflux_1/AU-Whr_WSA_ps.png}
\caption{Empirical Best Fit (BF) soil moisture PDF calibration for AU-Whr.}
\label{fig:ps_bf_AU-Whr}
\end{figure}
\begin{figure}[H]
\centering
\includegraphics[width=0.75\textwidth]{output_corrected/opt/figs/_ozflux_1/AU-Whr_WSA_ps.png}
\caption{Eco-Evolutionary Optimality (OPT) soil moisture PDF calibration for AU-Whr.}
\label{fig:ps_opt_AU-Whr}
\end{figure}
\newpage
\begin{figure}[H]
\centering
\includegraphics[width=0.9\textwidth]{output_corrected/bf/figs/_ozflux_1/AU-Whr_WSA_1.png}
\caption{MCMC diagnostics for derived parameters ($\Pi_T$, $\Pi_S$, $\Pi_R$, $\Pi_F$, $\epsilon$, and $AA$) under Empirical Best Fit (BF) for AU-Whr.}
\label{fig:1_bf_AU-Whr}
\end{figure}
\begin{figure}[H]
\centering
\includegraphics[width=0.9\textwidth]{output_corrected/opt/figs/_ozflux_1/AU-Whr_WSA_1.png}
\caption{MCMC diagnostics for derived parameters under Eco-Evolutionary Optimality (OPT) for AU-Whr.}
\label{fig:1_opt_AU-Whr}
\end{figure}
\newpage
\begin{figure}[H]
\centering
\includegraphics[width=0.9\textwidth]{output_corrected/bf/figs/_ozflux_1/AU-Whr_WSA_2.png}
\caption{MCMC diagnostics for calibrated traits ($P_{x50}$, $P_{g50}$, $k_{xl,max}$, $RAI$, $\beta_{ww}$, $s_{wilt}$, $s^*$, $\psi_{wilt}$, and $\psi^*$) under Empirical Best Fit (BF) for AU-Whr.}
\label{fig:2_bf_AU-Whr}
\end{figure}
\begin{figure}[H]
\centering
\includegraphics[width=0.9\textwidth]{output_corrected/opt/figs/_ozflux_1/AU-Whr_WSA_2.png}
\caption{MCMC diagnostics for calibrated traits under Eco-Evolutionary Optimality (OPT) for AU-Whr.}
\label{fig:2_opt_AU-Whr}
\end{figure}
\newpage
\subsection{AU-YarI --- Yanco Agricultural (GRA)}
Agricultural grassland site in New South Wales. Situated in a semi-arid zone, subject to seasonal cropping and pasture grazing.
\vskip 1.5em
\begin{figure}[H]
\centering
\includegraphics[width=0.75\textwidth]{output_corrected/bf/figs/_ozflux_1/AU-YarI_GRA_ps.png}
\caption{Empirical Best Fit (BF) soil moisture PDF calibration for AU-YarI.}
\label{fig:ps_bf_AU-YarI}
\end{figure}
\begin{figure}[H]
\centering
\includegraphics[width=0.75\textwidth]{output_corrected/opt/figs/_ozflux_1/AU-YarI_GRA_ps.png}
\caption{Eco-Evolutionary Optimality (OPT) soil moisture PDF calibration for AU-YarI.}
\label{fig:ps_opt_AU-YarI}
\end{figure}
\newpage
\begin{figure}[H]
\centering
\includegraphics[width=0.9\textwidth]{output_corrected/bf/figs/_ozflux_1/AU-YarI_GRA_1.png}
\caption{MCMC diagnostics for derived parameters ($\Pi_T$, $\Pi_S$, $\Pi_R$, $\Pi_F$, $\epsilon$, and $AA$) under Empirical Best Fit (BF) for AU-YarI.}
\label{fig:1_bf_AU-YarI}
\end{figure}
\begin{figure}[H]
\centering
\includegraphics[width=0.9\textwidth]{output_corrected/opt/figs/_ozflux_1/AU-YarI_GRA_1.png}
\caption{MCMC diagnostics for derived parameters under Eco-Evolutionary Optimality (OPT) for AU-YarI.}
\label{fig:1_opt_AU-YarI}
\end{figure}
\newpage
\begin{figure}[H]
\centering
\includegraphics[width=0.9\textwidth]{output_corrected/bf/figs/_ozflux_1/AU-YarI_GRA_2.png}
\caption{MCMC diagnostics for calibrated traits ($P_{x50}$, $P_{g50}$, $k_{xl,max}$, $RAI$, $\beta_{ww}$, $s_{wilt}$, $s^*$, $\psi_{wilt}$, and $\psi^*$) under Empirical Best Fit (BF) for AU-YarI.}
\label{fig:2_bf_AU-YarI}
\end{figure}
\begin{figure}[H]
\centering
\includegraphics[width=0.9\textwidth]{output_corrected/opt/figs/_ozflux_1/AU-YarI_GRA_2.png}
\caption{MCMC diagnostics for calibrated traits under Eco-Evolutionary Optimality (OPT) for AU-YarI.}
\label{fig:2_opt_AU-YarI}
\end{figure}
\newpage
\subsection{AU-Ync --- Yanco Grassland (GRA)}
Native grassland site in semi-arid New South Wales, characterized by high inter-annual climate variability driven by ENSO.
\vskip 1.5em
\begin{figure}[H]
\centering
\includegraphics[width=0.75\textwidth]{output_corrected/bf/figs/_ozflux_1/AU-Ync_GRA_ps.png}
\caption{Empirical Best Fit (BF) soil moisture PDF calibration for AU-Ync.}
\label{fig:ps_bf_AU-Ync}
\end{figure}
\begin{figure}[H]
\centering
\includegraphics[width=0.75\textwidth]{output_corrected/opt/figs/_ozflux_1/AU-Ync_GRA_ps.png}
\caption{Eco-Evolutionary Optimality (OPT) soil moisture PDF calibration for AU-Ync.}
\label{fig:ps_opt_AU-Ync}
\end{figure}
\newpage
\begin{figure}[H]
\centering
\includegraphics[width=0.9\textwidth]{output_corrected/bf/figs/_ozflux_1/AU-Ync_GRA_1.png}
\caption{MCMC diagnostics for derived parameters ($\Pi_T$, $\Pi_S$, $\Pi_R$, $\Pi_F$, $\epsilon$, and $AA$) under Empirical Best Fit (BF) for AU-Ync.}
\label{fig:1_bf_AU-Ync}
\end{figure}
\begin{figure}[H]
\centering
\includegraphics[width=0.9\textwidth]{output_corrected/opt/figs/_ozflux_1/AU-Ync_GRA_1.png}
\caption{MCMC diagnostics for derived parameters under Eco-Evolutionary Optimality (OPT) for AU-Ync.}
\label{fig:1_opt_AU-Ync}
\end{figure}
\newpage
\begin{figure}[H]
\centering
\includegraphics[width=0.9\textwidth]{output_corrected/bf/figs/_ozflux_1/AU-Ync_GRA_2.png}
\caption{MCMC diagnostics for calibrated traits ($P_{x50}$, $P_{g50}$, $k_{xl,max}$, $RAI$, $\beta_{ww}$, $s_{wilt}$, $s^*$, $\psi_{wilt}$, and $\psi^*$) under Empirical Best Fit (BF) for AU-Ync.}
\label{fig:2_bf_AU-Ync}
\end{figure}
\begin{figure}[H]
\centering
\includegraphics[width=0.9\textwidth]{output_corrected/opt/figs/_ozflux_1/AU-Ync_GRA_2.png}
\caption{MCMC diagnostics for calibrated traits under Eco-Evolutionary Optimality (OPT) for AU-Ync.}
\label{fig:2_opt_AU-Ync}
\end{figure}
\newpage

\end{document}
